\noindent\textbf{Author and source.} Lucas Mioranci, \emph{Algebraic hyperbolicity of very general hypersurfaces in homogeneous varieties}, exact arXiv:2307.10461v4 (18 March 2026); \'Epijournal de G\'eom\'etrie Alg\'ebrique 10 (2026), Article 5, DOI \href{https://doi.org/10.46298/epiga.2025.14608}{10.46298/epiga.2025.14608}.
\par\noindent\textbf{License.} Copyright the author. Source and this adapted presentation: Creative Commons Attribution--ShareAlike 4.0 International, \url{https://creativecommons.org/licenses/by-sa/4.0/}. No endorsement is implied.
\par\noindent\textbf{Scope and changes.} One target: the entire Theorem 1.4, restated as Theorem 4.3, with both degree ranges. The definition of algebraic hyperbolicity, Main Setting 1.3, and complete Sections 2--4 are retained verbatim from the fixed native source. The section numbering is restored; journal furniture is replaced by a standard wrapper. Section 5 applications and unrelated introductory discussion are not selected. No mathematical argument has been supplied or repaired.
\par\noindent\textbf{Dependency boundary.} Cited earlier results are allowed prerequisites; their statements and original references are retained, but their external proofs are not represented as reproduced. In particular this applies to the Coskun--Riedl normal-bundle and scroll tools. The retained author proofs include the new quantitative cases and the full Schubert argument. The final conjecture and illustrative examples are context, not counted results or asserted proofs.
\par\noindent\textbf{Preparation only.} No formal problem number, mathematical certification, or final global uniqueness is asserted. Later in-flight overlap after the saved 14 October 2026 cutoff is unconfirmed. AI assistance was used only for selection, metadata and mechanical packaging; no author AI-use assertion is inferred.
